\section*{Chapter \ref{ch:dv:greeks-and-the-hedging-pnl} --- Greeks and the Hedging P\&L}

\begin{solution}{exo:dv:greeks-and-the-hedging-pnl:1}
Delta 0.0230, gamma 0.1381, \omterm{def:dv:greeks-and-the-hedging-pnl:greeks}{vega} 0.2302 per point, \omterm{def:dv:greeks-and-the-hedging-pnl:greeks}{theta} $-0.0757$ per calendar day.
\end{solution}

\begin{solution}{exo:dv:greeks-and-the-hedging-pnl:2}
With $r=q=0$: $\Theta=-\tfrac12\sigma^2S^2\Gamma=-\tfrac12\times0.04\times10\,000\times0.1381
=-27.63$ per year, $-0.0757$ per day, as the analytic \omterm{def:dv:greeks-and-the-hedging-pnl:greeks}{theta}.
\end{solution}

\begin{solution}{exo:dv:greeks-and-the-hedging-pnl:3}
Gamma P\&L $5\,000\times0.02^2=2.00$; \omterm{def:dv:greeks-and-the-hedging-pnl:greeks}{theta} $-5\,000\times0.2^2/365=-0.55$; net $+1.45$.
\end{solution}

\begin{solution}{exo:dv:greeks-and-the-hedging-pnl:4}
Each interval pays $\tfrac12\Gamma S^2(x^2-\sigma^2\delta t)$, quadratic in the return
$x$. A share that drifts to 120 in small steps (each below the break-even move) loses
\omterm{def:dv:greeks-and-the-hedging-pnl:greeks}{theta} every day: the long \omterm{def:dv:greeks-and-the-hedging-pnl:cashgamma}{straddle}, hedged, loses money although unhedged it would
have paid 20 at expiry. The hedge sold the drift away as it happened.
\end{solution}

\begin{solution}{exo:dv:greeks-and-the-hedging-pnl:5}
By four: the error falls like $1/\sqrt N$. Each trade then rebalances a delta change
about half as large, so the total traded quantity, and the cost, grow like $\sqrt N$:
twice the cost for half the noise (the Leland trade-off).
\end{solution}

\begin{solution}{exo:dv:greeks-and-the-hedging-pnl:6}
$\sqrt{\pi/4}\times23.02\times0.2/\sqrt{21}=0.890$, within 3\% of the simulated 0.865.
\end{solution}

\begin{solution}{exo:dv:greeks-and-the-hedging-pnl:7}
Mean $-1.147$, standard deviation 1.082: the mean is unchanged, and with daily hedges
the discrete-hedging noise, not the choice of volatility, dominates the spread.
\end{solution}

\begin{solution}{exo:dv:greeks-and-the-hedging-pnl:8}
A short \omterm{def:dv:greeks-and-the-hedging-pnl:cashgamma}{straddle} collects \omterm{def:dv:greeks-and-the-hedging-pnl:greeks}{theta} on most days and pays gamma on the few days with large
moves; with realised above implied, the quadratic losses on large-move days exceed the
steady gains. A high share of winning days is the signature of selling
optionality, not evidence of a pricing error; the loss is the gamma--theta identity at
work.
\end{solution}

\begin{solution}{pb:dv:greeks-and-the-hedging-pnl:1}
\textbf{1.} 4.6059 per \omterm{def:dv:greeks-and-the-hedging-pnl:cashgamma}{straddle}; USD~460\,595.
\textbf{2.} Each \omterm{def:dv:greeks-and-the-hedging-pnl:cashgamma}{straddle} has delta 0.0230; short 1\,000 with multiplier 100 is $-2\,303$
shares; the desk buys 2\,303 shares.
\textbf{3.} \omterm{def:dv:greeks-and-the-hedging-pnl:cashgamma}{Cash gamma} 690.70; gamma per 1\% move 13.81 (of cash delta).
\textbf{4.} \omterm{def:dv:greeks-and-the-hedging-pnl:greeks}{Vega} USD~23\,023 per point (short, so a loss when volatility rises); \omterm{def:dv:greeks-and-the-hedging-pnl:greeks}{theta}
USD~7\,569 a day in the desk's favour.
\textbf{5.} 1.047 over 365 days; 1.260 over 252.
\textbf{6.} 5.7570.
\textbf{7.} 1.1510 per \omterm{def:dv:greeks-and-the-hedging-pnl:cashgamma}{straddle}; USD~115\,104.
\textbf{8.} $5\times0.2302=1.1512$.
\textbf{9.} The \omterm{def:dv:greeks-and-the-hedging-pnl:cashgamma}{straddle}'s value is nearly linear in volatility near the money (\omterm{def:dv:greeks-and-the-hedging-pnl:greeks}{volga} is
small), so \omterm{def:dv:greeks-and-the-hedging-pnl:greeks}{vega} times the gap is an accurate first-order estimate.
\textbf{10.} No: continuous hedging gives the same expectation, with no dispersion
coming from the hedging schedule.
\textbf{11.} Mean $-1.145$, standard deviation 1.183.
\textbf{12.} 0.620.
\textbf{13.} 87.3\%.
\textbf{14.} Mean $-1.147$, standard deviation 1.082.
\textbf{15.} Mean 0.003, standard deviation 0.865: pure discrete-hedging noise.
\textbf{16.} That the share would realise at most 20 over the month, after costs and
with a margin for the hedging noise.
\textbf{17.} Limit the \omterm{def:dv:greeks-and-the-hedging-pnl:greeks}{vega} (and the gamma near expiry) so that a plausible volatility
shock, say ten points, and two standard deviations of hedging noise stay within the
monthly loss limit.
\textbf{18.} The \omterm{def:dv:greeks-and-the-hedging-pnl:greeks}{vega}, read with a scenario for realised volatility, and the gamma
concentration near the strike.
\textbf{19.} An expected loss of USD~115\,100 and a standard deviation of USD~118\,300
with daily hedging.
\textbf{20.} The \omterm{def:dv:greeks-and-the-hedging-pnl:cashgamma}{cash gamma} times the difference between realised and implied variance,
summed over the life of the option.
\end{solution}

\begin{solution}{iq:dv:greeks-and-the-hedging-pnl:1}
It lost a week of \omterm{def:dv:greeks-and-the-hedging-pnl:greeks}{theta}: the option decayed and no gamma P\&L came in to pay for it.
By the gamma--theta identity, a hedged long option earns only when the underlying moves
more than the implied volatility assumes.
\iqlookfor{\omterm{def:dv:greeks-and-the-hedging-pnl:greeks}{theta} as the rent for gamma.}
\end{solution}

\begin{solution}{iq:dv:greeks-and-the-hedging-pnl:2}
$\Theta\approx-\tfrac12\Gamma S^2\sigma^2$: the time decay of a hedged option is
exactly what its gamma earns if the underlying moves by one standard deviation of the
implied volatility each day. Being long gamma means paying \omterm{def:dv:greeks-and-the-hedging-pnl:greeks}{theta} for convexity.
\iqlookfor{the identity, and ``break-even at implied''.}
\end{solution}

\begin{solution}{iq:dv:greeks-and-the-hedging-pnl:3}
Taylor-expand the option: $\Theta\delta t+\Delta\delta S+\tfrac12\Gamma\delta S^2$; the
hedge removes $\Delta\delta S$; the identity replaces $\Theta$ by $-\tfrac12\Gamma S^2
\sigma_{\mathrm{imp}}^2$ (zero rates): P\&L $=\tfrac12\Gamma S^2\bigl((\delta S/S)^2-
\sigma_{\mathrm{imp}}^2\delta t\bigr)$.
\iqlookfor{the second-order expansion and the cancellation.}
\end{solution}

\begin{solution}{iq:dv:greeks-and-the-hedging-pnl:4}
At your forecast, if it is right, the P\&L at expiry is locked in, but the daily
mark-to-market swings; at implied, daily P\&L is explained by gamma and \omterm{def:dv:greeks-and-the-hedging-pnl:greeks}{theta}, but the
total depends on where the large moves happen relative to the strike. Hedging at
implied is the default because the forecast is uncertain and because risk reports
and limits are in implied terms.
\iqlookfor{the two properties, and why desks choose implied.}
\end{solution}

\begin{solution}{iq:dv:greeks-and-the-hedging-pnl:5}
Monte Carlo: the up and down repricings use independent random numbers, so their
difference is noise divided by $h^2$; use common random numbers, pathwise or
likelihood-ratio estimators (chapter 23). Digital near expiry: the payoff is
discontinuous, gamma is a spike narrower than the bump, and the difference measures
the bump size; report gamma over a bump that matches the risk horizon, and manage
the digital by a call spread (chapter 15).
\iqlookfor{noise from independent draws, and non-smooth payoffs.}
\end{solution}

\begin{solution}{iq:dv:greeks-and-the-hedging-pnl:6}
\omterm{def:dv:greeks-and-the-hedging-pnl:greeks}{Vega} sums gamma over the life of the options, but the P\&L weights each day's realised
variance by that day's \omterm{def:dv:greeks-and-the-hedging-pnl:cashgamma}{cash gamma}. A book whose gamma sat at strikes the spot visited
on its volatile days earns more than one whose gamma sat elsewhere; the difference can
be of the order of the whole expected P\&L, as the dispersion of
\cref{fig:dv:greeks-and-the-hedging-pnl:which} shows.
\iqlookfor{gamma weighting in time and in spot.}
\end{solution}
